\documentclass[aspectratio=169, 11pt]{beamer}

\usepackage[utf8]{inputenc}
\usepackage[T1]{fontenc}
\usepackage[ngerman]{babel}
\usepackage{helvet}
\renewcommand{\familydefault}{\sfdefault}
\usepackage{xcolor}

\usetheme[progressbar=none, block=fill, numbering=fraction]{metropolis}

\definecolor{darkslate}{HTML}{0F172A}
\definecolor{slategray}{HTML}{475569}
\definecolor{lightbg}{HTML}{F8FAFC}
\definecolor{rulecolor}{HTML}{CBD5E1}
\definecolor{mutedtext}{HTML}{64748B}

\setbeamercolor{normal text}{fg=darkslate, bg=white}
\setbeamercolor{alerted text}{fg=slategray}
\setbeamercolor{frametitle}{fg=white, bg=darkslate}
\setbeamercolor{block title}{fg=white, bg=darkslate}
\setbeamercolor{block body}{fg=darkslate, bg=lightbg}
\setbeamercolor{title separator}{fg=slategray}
\setbeamercolor{title}{fg=white}
\setbeamercolor{subtitle}{fg=white!70!darkslate}
\setbeamercolor{author}{fg=white!70!darkslate}
\setbeamercolor{date}{fg=white!60!darkslate}
\setbeamercolor{background canvas}{bg=white}

\setbeamerfont{title}{size=\large, series=\bfseries}
\setbeamerfont{subtitle}{size=\normalsize, series=\mdseries}
\setbeamerfont{frametitle}{size=\normalsize, series=\bfseries}
\setbeamerfont{author}{size=\small}
\setbeamerfont{block title}{size=\small, series=\bfseries}
\setbeamerfont{block body}{size=\small}

\setbeamertemplate{itemize item}{\small\raise1pt\hbox{--}}
\setbeamertemplate{itemize subitem}{\scriptsize\raise1pt\hbox{--}}

\newcommand{\hinweis}[1]{{\scriptsize\color{mutedtext}\textit{#1}}}

% -------------------------------------------------------
\title{Bottleneck-Analyse im Karosseriebau}
\subtitle{Zwischenstand}
\author{Vorname Name \quad Vorname Name}
\institute{Datenanalyse im Produktionsumfeld}
\date{Sommersemester 2025}

% -------------------------------------------------------
\begin{document}

\maketitle

% ----------------------------------------------------------------
\begin{frame}{Was haben wir untersucht -- und was haben wir gefunden?}

  \begin{columns}[T, onlytextwidth]
    \begin{column}{0.55\textwidth}

      % \includegraphics[width=\linewidth]{plots/eür_plot.png}
      \begin{center}
        \textcolor{rulecolor}{\rule{\linewidth}{0.4pt}}
        \vspace{1.5cm}
        {\small\textcolor{mutedtext}{Grafik einfügen}}
        \vspace{1.5cm}
        \textcolor{rulecolor}{\rule{\linewidth}{0.4pt}}
      \end{center}

    \end{column}

    \begin{column}{0.41\textwidth}

      \textbf{Unsere Frage}\\[4pt]
      {\small\textit{[Eine Frage in einem Satz]}}

      \vspace{10pt}
      \textbf{Kernergebnis}\\[4pt]
      {\small\textit{[Was zeigt der Plot? Konkret:
      welche Station, welcher Effekt, welche Grössenordnung?]}}

      \vspace{10pt}
      \textbf{Vorgehen}\\[4pt]
      {\small\textit{[2--3 Sätze: Welche Daten,
      welche Methode, warum so?]}}

    \end{column}
  \end{columns}

  \hinweis{Eine Grafik, eine Aussage. Nicht mehr.}

\end{frame}

% ----------------------------------------------------------------
\begin{frame}{Welche Fragen haben wir jetzt?}

  \begin{columns}[T, onlytextwidth]
    \begin{column}{0.47\textwidth}

      \textbf{Fragen, die wir selbst beantworten können}\\[6pt]
      \begin{itemize}
        \item \textit{[Folgefrage, die mit den vorhandenen Daten\\
          beantwortbar ist -- nächster Analyseschritt]}
        \item \textit{[Weitere Folgefrage]}
      \end{itemize}

    \end{column}

    \begin{column}{0.47\textwidth}

      \textbf{Fragen, für die wir euch brauchen}\\[6pt]
      \begin{itemize}
        \item \textit{[Was fehlt uns an Kontext oder Daten?
          Z.\,B.: ,,Was passiert konkret an Station X?``]}
        \item \textit{[Weitere Frage an den Konzern]}
      \end{itemize}

    \end{column}
  \end{columns}

  \vspace{14pt}

  \begin{block}{Unser nächster Schritt}
    \textit{[Was analysiert ihr als nächstes? Ein konkreter Satz.]}
  \end{block}

  \hinweis{Die rechte Spalte ist genauso wichtig wie die linke.
  Zeigt, dass ihr wisst was ihr noch nicht wisst.}

\end{frame}

\end{document}
